%!TEX root = ./../../springer-mono-format/book.tex
%!TEX root = ./../../thesis-trier.tex

Here, we introduce morphisms, endomorphisms, isomorphisms, and automorphisms
of automata and collect basic facts about them.

\subsection{Definitions and Basic Facts}

Let $\semiautomat{A}{Q}{\Sigma}{\delta}$ and $\semiautomat{B}{S}{\Sigma}{\mu}$
be two semiautomata. 

\IN{morphism!of (semi-)automata}
\IN{endomorphism!semiautomaton}
\IN{isomorphism!semiautomaton}
\IN{automorphism!semiautomaton}
\IN{isomorphic automata}
\begin{enumerate}
	\item A \emph{(semiautomaton) (homo-)morphism} from $\mathcal A$ to $\mathcal B$
	(or between $\mathcal A$ and $\mathcal B$) is a function
	$\varphi \colon Q \to S$
	such that for each $q \in Q$ and $x \in \Sigma$ we have
	\[
	\applyfunc{\applytrans{\delta}{q}{x}}{\varphi}
	=
	\applytrans{\mu}{\applyfunc{q}{\varphi}}{x}.
	\]
	Equivalently, $\varphi$ is a morphism if and only if
	$\concat{\delta_x}{\varphi} = \concat{\varphi}{\mu_x}$
	for all $x \in \Sigma$.
	Thus, regardless of which of the two paths we follow in the diagram below,
	we end up at the same state.
	\[
	\begin{tikzcd}[row sep=1cm,column sep=1.25cm]
		Q \ar{r}{\delta(\cdot, x)} \ar[swap]{d}{\varphi}
		& Q \ar{d}{\varphi} \\
		S \ar[swap]{r}{\mu(\cdot, x)}
		& S
	\end{tikzcd}
	\]
	If $\varphi$ is surjective, we say that it maps $\mathcal A$ \emph{onto}
	$\mathcal B$.
	
	\item  A \emph{(semiautomaton) endomorphism}
	is a morphism with $\mathcal A = \mathcal B$.
	
	\item A \emph{(semiautomaton) isomorphism} is a morphism $\varphi \colon Q \to S$ such that
	there exists a morphism $\psi \colon S \to Q$ with $\concat{\varphi}{\psi} = \operatorname{id}_Q$
	and $\concat{\psi}{\varphi} = \operatorname{id}_S$.
	The automata~$\mathcal A$ 
	and~$\mathcal B$ are called \emph{isomorphic}, written \[ \mathcal A \cong \mathcal B, \] \label{def:isomorphic}
	if there exists an isomorphism between them.
	
	\item A \emph{(semiautomaton) automorphism} is an endomorphism that is an isomorphism.
\end{enumerate}

The set of endomorphisms of a given automaton forms a transformation semigroup
and the set of automorphisms forms a permutation group.

\begin{example} \label{prelim:aut:ex:mappings}
	\label{prelim:ex:SCC} \label{prelim:ex:SCC2}
	Consider the automata from \cref{prelim:aut:fig:mappings}, where the dashed lines indicate a 
	mapping of their state sets.
		
	In the automaton from \cref{prelim:aut:fig:mappingmorph}
	the mapping is an endomorphism that does not map the strongly connected
	component $\{q_0\}$ to a strongly connected component, but to $\{q_1\}$.	
	Another endomorphism of the same automaton is given in \cref{prelim:aut:fig:mappingmorph0}.
	Together with the identity transformation, this automaton has exactly three endomorphisms.
	
	In \cref{prelim:aut:fig:mappingnotmorph0,prelim:aut:fig:mappingotmorph},
	the mapping $\varphi$ is not a morphism since in both cases we have
	$\varphi(\delta(q_0, a)) = q_1 \ne q_2 = \delta(\varphi(q_0), a)$.
	
	\begin{figure}[ht]
	\centering
	\begin{subfigure}{0.25\textwidth} \centering
		\begin{tikzpicture}[node distance = 2cm, on grid, auto, ->]
			\node (q0) [state] {$q_0$};
			\node (q1) [state, below of = q0] {$q_1$};
			\node (q2) [state, below of = q1] {$q_2$};
			
			\path (q0) edge node {$a$} (q1) 
			(q1) edge [bend left] node {$a$} (q2)
			(q2) edge [bend left] node {$a$} (q1);
			
			\path (q0) edge [dashed,bend left=2cm] (q1);
			\path (q1) edge [dashed, bend left=3cm] (q2)
			(q2) edge [dashed, bend left=3cm] (q1);
		\end{tikzpicture}
		\caption{An endomorphism.}
		\label{prelim:aut:fig:mappingmorph}
	\end{subfigure}%
	\hfill
	\begin{subfigure}{0.25\textwidth} \centering
		\begin{tikzpicture}[node distance = 2cm, on grid, auto, ->]
			\node (q0) [state] {$q_0$};
			\node (q1) [state, below of = q0] {$q_1$};
			\node (q2) [state, below of = q1] {$q_2$};
			
			\path (q0) edge node {$a$} (q1) 
			(q1) edge [bend left] node {$a$} (q2)
			(q2) edge [bend left] node {$a$} (q1);
			
			\path (q0) edge [dashed,bend left=2cm] (q2);
			\path (q1) edge [dashed, loop left] (q1)
			(q2) edge [dashed, loop left] (q2);
		\end{tikzpicture}
		\caption{An endomorphism.}
		\label{prelim:aut:fig:mappingmorph0}
	\end{subfigure}%
	\hfill
	\begin{subfigure}{0.25\textwidth} \centering
		\begin{tikzpicture}[node distance = 2cm, on grid, auto, ->]
			\node (q0) [state] {$q_0$};
			\node (q1) [state, below of = q0] {$q_1$};
			\node (q2) [state, below of = q1] {$q_2$};
			
			\path (q0) edge node {$a$} (q1) 
			(q1) edge [bend left] node {$a$} (q2)
			(q2) edge [bend left] node {$a$} (q1);
			
			\path (q0) edge [dashed,bend left=2cm] (q1);
			\path (q1) edge [dashed, loop left] (q1)
			(q2) edge [dashed, loop left] (q2);
		\end{tikzpicture}
		\caption{Not a morphism.}
		\label{prelim:aut:fig:mappingnotmorph0}
	\end{subfigure}%
	\hfill
	\begin{subfigure}{0.25\textwidth} \centering
		\begin{tikzpicture}[node distance = 2cm, on grid, auto, ->]
			\node (q0) [state] {$q_0$};
			\node (q1) [state, below of = q0] {$q_1$};
			\node (q2) [state, below of = q1] {$q_2$};
			
			\path (q0) edge [loop right ]node {$a$} (q0) 
			(q1) edge [bend left] node {$a$} (q2)
			(q2) edge [bend left] node {$a$} (q1);
			
			\path (q0) edge [dashed,bend left=2cm] (q1);
			\path (q1) edge [dashed, bend left=3cm] (q2)
			(q2) edge [dashed, bend left=3cm] (q1);
		\end{tikzpicture}
		\caption{Not a morphism.}
		\label{prelim:aut:fig:mappingotmorph}
	\end{subfigure}
	\caption[Example morphisms and non-morphisms on automata]{Automata and mappings of their states indicated by dashed lines, see \cref{prelim:aut:ex:mappings}.}
	\label{prelim:aut:fig:mappings}
\end{figure}
\end{example}

\begin{example} \label{prelim:aut:ex:autnomorph0}
	Let $\semiautomat{A}{\{q_1, q_2\}}{\{a_1,a_2\}}{\delta}$
	with $\delta(q_1, a_i) = \delta(q_2, a_i) = q_i$
	for $i \in \{1,2\}$, see \cref{prelim:aut:fig:autnomorph}.
	If~$\varphi \colon Q \to Q$ is an endomorphism,
	then $\varphi(q_i) = \varphi(\delta(q_i, a_i)) = \delta(\varphi(q_i), a_i) = q_i$
	for $i \in \{1,2\}$.
	Thus the identity transformation is the only endomorphism of~$\mathcal A$.	
	
	Later, in \cref{prelim:automata:morphism:cor:endo-sync-SCC}, we will show that, in general, a 
	strongly connected and synchronizing automaton only admits the identity transformation as an endomorphism.
	
	\begin{figure}[ht]
		\centering	
		\begin{tikzpicture}[node distance = 2cm, on grid, auto, ->]
			\node (q0) [state] {$q_1$};
			\node (q1) [state, right of = q0] {$q_2$};
			
			\path (q0) edge [bend left=10] node {$a_2$} (q1) 
			(q1) edge [loop right] node {$a_2$} (q1);
			
			\path (q1) edge [bend left=10] node {$a_1$} (q0)
			(q0) edge [loop left] node {$a_1$} (q0);
		\end{tikzpicture}
		\caption[Automaton not admitting any non-identity endomorphism.]{Automaton not admitting any non-identity endomorphism, see \cref{prelim:aut:ex:autnomorph0}.}
		\label{prelim:aut:fig:autnomorph}
	\end{figure}
\end{example}

Next, we give some alternative characterizations for morphisms.

\begin{proposition} \label{prelim:automata:morphism:prop:principal-upper-set}
	Let $\varphi \colon Q \to S$ be a function. Then the following are equivalent.
	
	\begin{enumerate}
		\item \label{prelim:automata:morphism:prop:principal-upper-set:1} $\varphi$ is a morphism.
		
		\item \label{prelim:automata:morphism:prop:principal-upper-set:2} For all $p \in Q$ and all $u \in \Sigma^*$
		we have
		\[
		\varphi(\delta(p, u)) = \mu(\varphi(p), u)
		\]
		
		\item \label{prelim:automata:morphism:prop:principal-upper-set:3} There exists $T \subseteq Q$ with $Q = \bigcup_{p \in T} \upperset{\{p\}}$
		such that for all $p \in T$ and all $u \in \Sigma^*$ we have
		\[
		\varphi(\delta(p, u)) = \mu(\varphi(p), u)
		\]
	\end{enumerate}
\end{proposition}
\begin{proof}
	$\ref{prelim:automata:morphism:prop:principal-upper-set:1}\Rightarrow\ref{prelim:automata:morphism:prop:principal-upper-set:2}$: This clearly holds for $u = \varepsilon$.
	Next, let $u = va$ for some $v \in \Sigma^*$ and $a \in \Sigma$.
	Inductively, we have $\varphi(\delta(p, v)) = \mu(\varphi(p), v)$.
	Using that~$\varphi$ is a morphism, this gives $\varphi(\delta(p, u)) 
	= \varphi(\delta(\delta(p, v), a))
	= \mu(\varphi(\delta(p, v)), a)
	= \mu(\mu(\varphi(p), v), a)
	= \mu(\varphi(p), u)$.
	Then the claim follows by induction.
	
	$\ref{prelim:automata:morphism:prop:principal-upper-set:2}\Rightarrow\ref{prelim:automata:morphism:prop:principal-upper-set:3}$: By assumption, this holds for $T = Q$.
	
	$\ref{prelim:automata:morphism:prop:principal-upper-set:3}\Rightarrow\ref{prelim:automata:morphism:prop:principal-upper-set:1}$: Let $a \in \Sigma$ and $q \in Q$.
	By assumption, there exists $u \in \Sigma^*$ such that $q = \delta(p, u)$
	for some $p \in T$. Set $s = \varphi(p)$.
	Then we have
	\begin{align*}
		\varphi(\delta(q, a))
		& = \varphi(\delta(p, ua)) \tag{$q = \delta(p,u)$} \\
		& = \mu(s, ua)\tag{$\varphi$ morphism}  \\
		& = \mu(\mu(s, u), a) \tag{Def.~transition function} \\
		& = \mu(\varphi(\delta(p, u)), a) \tag{$\varphi(\delta(p, u)) = \mu(s, u)$} \\
		& = \mu(\varphi(q), a). \tag{$q = \delta(p,u)$}
	\end{align*}
	Thus, $\varphi$ is a morphism.
\end{proof}

By \cref{prelim:automata:morphism:prop:principal-upper-set:3} of \cref{prelim:automata:morphism:prop:principal-upper-set}, a morphism
is completely specified by the images of $\varphi(p_1), \ldots, \varphi(p_m)$
for some $p_1, \ldots, p_m \in Q$ such that 
$Q = \upperset{\{p_1\}} \cup \cdots \cup \upperset{\{p_m\}}$.

\begin{corollary}
	If~$\mathcal A$ is cyclic, there are at most~$|S|$ morphisms
	from~$\mathcal A$ into~$\mathcal B$.
	In particular, there are at most~$|Q|$ endomorphisms of~$\mathcal A$.
\end{corollary}
\begin{proof}
	By \cref{prelim:automata:morphism:cor:two-morp-upper-set} 
	two morphisms having the same image for a state~$p$ with $Q = \upperset{\{p\}}$
	are equal. Since we have at most $|S|$ distinct images for this state,
	the claim follows.
\end{proof}

\begin{example}
	The might be strictly less than $|Q|$ endomorphisms in a cyclic automaton.
	For example, the automaton from \cref{prelim:aut:ex:autnomorph0}
	is cyclic and admits no non-identity morphism.
\end{example}

\IN{morphic image}
In the rest of this section, let  $\varphi \colon Q \to S$ 
be a morphism. For $T \subseteq Q$, we call $\applyfunc{T}{\varphi}$ the \emph{morphic image (of $T$)}.

By our next statement, isomorphisms are precisely the bijective morphisms. 
Note that a similar statement does not hold for order morphisms, see \cref{prelim:ex:order:bijmorph}.

\begin{proposition}\label{prelim:prop:iso-bijmorph}
	The morphism $\varphi \colon Q \to S$ is an isomorphism if and only if it is bijective.
\end{proposition}
\begin{proof}
	By \cref{prelim:prop:surj-inj-inverses}, existence of a map $\psi \colon S \to Q$ 
	such that $\concat{\varphi}{\psi} = \operatorname{id}_Q$
	and $\concat{\psi}{\varphi} = \operatorname{id}_S$
	is equivalent to bijectivity of $\varphi$. So, an isomorphism is bijective.
	
	Next, assume~$\varphi$ is bijective. 
	Since~$\varphi$ is surjective,
	we have $\concat{\psi}{\varphi} = \operatorname{id}_S$
	for some $\psi \colon S \to Q$ 
	by \cref{prelim:prop:surj-inj-inverses}.
	Let $s \in S$ and $x \in \Sigma$.
	Then we find
	\begin{align*}
		\varphi(\delta(\psi(s), x))
		& = \mu(\varphi(\psi(s)), x) \tag{$\varphi$ morphism} \\
		& = \mu(s, x) \tag{$\concat{\psi}{\varphi} = \operatorname{id}_S$} \\
		& = \varphi(\psi(\mu(s, x))). \tag{$\concat{\psi}{\varphi} = \operatorname{id}_S$}
	\end{align*}
	Since~$\varphi$ is injective, $\varphi(\delta(\psi(s), x)) = 
	\varphi(\psi(\mu(s, x)))$
	implies $\delta(\psi(s), x) = \psi(\mu(s, x))$.
	Thus $\psi$ is a morphism.
\end{proof}

\begin{remark}
	By \cref{prelim:prop:surj-inj-inverses} 
	a morphism $\varphi \colon Q \to S$
	is injective (resp.~surjective) if
	there exists a morphim $\psi \colon S \to Q$
	such that $\concat{\varphi}{\psi}$ (resp.~$\concat{\psi}{\varphi}$).
	However, unlike for functions,
	the converse implication does not hold in general.	

	To show this, for $i \in \{1,2,3\}$, let $\mathcal A_i = (Q_i, \{a\}, \delta_i)$
	where $Q_1 = \{ q_1, q_2 \}$, $Q_2 = \{ q_1, q_2, t \}$, $Q_3 = \{t\}$,
	and where the letter~$a$ swaps the states $q_1$ and~$q_2$ in $\mathcal A_1$ and~$\mathcal A_2$
	and fixes the state~$t$ in $\mathcal A_2$ and~$\mathcal A_3$.	
	Define $\varphi \colon Q_1 \to Q_2$ by $\varphi(q_1)=q_1$, $\varphi(q_2)=q_2$,
	and $\psi \colon Q_1 \to Q_3$ by $\psi(q_1)=\psi(q_2)=t$.
	Then $\varphi$ is injective and $\psi$ is surjective
	and both maps are morphisms.	
	However, there exists no morphism from $\mathcal A_2$ to $\mathcal A_1$,
	and no morphism from $\mathcal A_3$ to $\mathcal A_1$.
\end{remark}

\begin{proposition} \label{prelim:aut:prop:morph:trap-to-trap}
	If~$t \in Q$ is a trap state, then $\varphi(t)$ is a trap state.
\end{proposition}
\begin{proof}
	Let $x \in \Sigma$.
	Then $\mu(\varphi(t), x) = \varphi(\delta(t,x)) = \varphi(t)$
	and so~$\varphi(t)$ is a trap state.
\end{proof}

\begin{proposition} \label{prelim:automata:morphism:prop:surj-morphism-preserves-sync}
	If~$\varphi$ is surjective, the image of a synchronizing state is a synchronizing state.
\end{proposition} 
\begin{proof}
	Let $\delta(Q, u) = \{q\}$ for $u \in \Sigma^*$ and $q \in Q$.
	Then $\mu(S, u) = \mu(\varphi(Q), u) = \varphi(\delta(Q,u)) = \{\varphi(q)\}$
\end{proof}

Next, a morphism of automata induces a morphism of the underlying digraphs.

\begin{proposition} \label{prelim:aut:prop:morph-digraph}
	$\varphi \colon Q \to S$ is also a graph morphism
	of the automaton digraphs of~$\mathcal A$ and~$\mathcal B$.
\end{proposition}
\begin{proof}
	For $p \in Q$ and $x \in \Sigma$, the edge $(p, \delta(p,x))$
	is mapped to the edge $(\varphi(p), \mu(\varphi(p), x))$
	by~$\varphi$ in the respective automaton digraphs,
	which gives the claim.
\end{proof}

The following two examples gives mappings that are graph morphisms of the automaton digraphs,
but not automaton morphisms.

\begin{example}
Let $\semiautomat{A}{\{p,q\}}{\{a,b\}}{\delta}$
with $\delta(p, x) = q$ for $x \in \{a,b\}$, $\delta(q, a) = p$, and $\delta(q, b) = q$.
Then the transformation~$\varphi$ that maps all states to~$q$
is a graph morphism of the automaton digraph of~$\mathcal A$,
but it is not an endomorphism
since $\delta(\varphi(q), a) = p$ and $\varphi(\delta(q, a)) = q$.
\end{example}

\begin{example}
Let $\semiautomat{A}{\{q_1,\ldots, q_n\}}{\Sigma}{\delta}$ with $n \ge 2$
be such that for each pair of distinct states $q_i, q_j$ there exists 
a letter $a_{ij}$ inducing a simple idempotent mapping $q_i$ to~$q_j$.
Every permutation of $\{q_1, \ldots, q_n\}$
is a graph automorphism of the automaton digraph of~$\mathcal A$.
However, $\mathcal A$ has no non-identity
automaton morphism
since if $\varphi(q_i) = q_j$ with $i \ne j$,
then $\varphi(\delta(q_i, a_{ji})) = q_j$,
but $\delta(\varphi(q_i), a_{ji}) = q_i$.
\end{example}



\subsection{Morphisms and the Reachability Order}

In this section, let $\mathcal A = (Q, \Sigma, \delta)$ and $\mathcal B = (S, \Sigma, \mu)$
be two automata.

First, we show that two morphisms agree on an upper set of the form $\upperset{\{p\}}$
if and only if they agree at~$p$.

\begin{proposition} \label{prelim:automata:morphism:cor:two-morp-upper-set}
For morphisms $\varphi, \psi \colon Q \to S$ and $p \in Q$, we have
  \begin{equation}
	\varphi(p) = \psi(p) \Rightarrow \forall q \in \upperset{\{p\}} \colon \varphi(q) = \psi(q) \label{prelim:eqn:two-morph-equal} 
  \end{equation}
 In particular, two morphism $\varphi, \psi$ agree on a strongly connected component
 or have different images for each state of the strongly connected component.
\end{proposition}
\begin{proof}
	Let $q \in \upperset{\{p\}}$.
	Then we have $q = \delta(p, u)$ for some $u \in \Sigma^*$.
	By \cref{prelim:automata:morphism:prop:principal-upper-set}
	we find $\varphi(\delta(p, u)) = \mu(\varphi(p), u) =  \mu(\psi(p), u) = \psi(\delta(p, u))$.
\end{proof}

For the remainder of this section, suppose $\varphi \colon Q \to S$ is a morphism.
Let $p,q \in Q$. 
Then $\varphi$ is an order morphism with respect to~$\lereach$.
Indeed, if $p \lereach q$, then $q = \delta(p,u)$ for some $u \in \Sigma^*$.
Consequently, 
$
\varphi(q) = \varphi(\delta(p,u)) = \mu(\varphi(p),u),
$
so $\varphi(p) \lereach \varphi(q)$.
Thus, we have shown that\footnote{This also 
	follows from \cref{prelim:aut:prop:morph-digraph}
	since a morphism maps paths to paths in the automaton digraphs.}
\begin{equation}
	p \lereach q \Rightarrow \varphi(p) \lereach \varphi(q). \label{prelim:eqn:morph-reachability-order}	
\end{equation}

\begin{proposition} \label{prelim:automata:morpisms:reachability:prop:images}
Let $\varphi \colon Q \to S$ be a morphism. %Then we have the following:

\begin{enumerate}
\item \label{prelim:automata:morpisms:reachability:prop:images:principal-upper-set} Let $T \subseteq Q$.
Then we have
\begin{equation}  \label{prelim:eqn:morph-principal-upper-set}
\varphi(\upperset{T}) = \upperset{\varphi(T)}
\end{equation}

\item  \label{prelim:fact:morphism-sink} \label{prelim:automata:morpisms:reachability:prop:images:preservation}
$\varphi(T)$ is an upper set (resp.~unilaterally connected, strongly connected, a sink component) for an upper set (resp.~unilaterally connected, strongly connected, sink component) $T \subseteq Q$.

\item \label{prelim:automata:morpisms:reachability:prop:images:twoSCC} For two strongly connected components $T \subseteq Q$ and $R \subseteq S$,
 we have $\varphi(T) \cap R = \emptyset$ or $\varphi(T) \subseteq R$.
 
\item \label{prelim:automata:morpisms:reachability:prop:images:inverse} Let $T \subseteq Q$
 be an upper set. Then $\varphi^{-1}(T)$ is also an upper set.
\end{enumerate}
\end{proposition}
\begin{proof}
\ref{prelim:automata:morpisms:reachability:prop:images:principal-upper-set}: 
We have
\begin{align*}
	\varphi(\upperset{T}) 
	& = \varphi(\{ \delta(p, u) \mid p \in T \land u \in \Sigma^* \}) \\
	& = \{ \varphi(\delta(p, u)) \mid p \in T \land u \in \Sigma^* \} \\
	& = \{ \mu(\varphi(p), u) \mid p \in T \land u \in \Sigma^* \} \\
	& = \upperset{\varphi(T)}.
\end{align*}

\ref{prelim:automata:morpisms:reachability:prop:images:preservation}: Let $T \subseteq Q$ be an upper set. 
Then $T = \upperset{T}$
and, with \cref{prelim:automata:morpisms:reachability:prop:images:principal-upper-set}, we find 
$\varphi(T) = \varphi(\upperset{T}) = \upperset{\varphi(T)}$.
Since $\upperset{\varphi(T)} \subseteq \varphi(T)$, $\varphi(T)$ is an upper set.

For unilaterally connected~$T$, consider $\varphi(p), \varphi(q)$ for $p,q \in T$.
Since $p \lereach q$ or $q \lereach p$, we have $\varphi(p) \lereach \varphi(q)$
or $\varphi(q) \lereach \varphi(p)$ and so~$\varphi(T)$
is unilaterally connected.

By \cref{prelim:prop:morphism_equiv_class}, $\varphi(T)$ is strongly connected
if~$T$ is strongly connected.

Lastly, suppose~$T$ is a sink component.
Then $\varphi(T)$ is strongly connected and an upper set
since both properties are preserved as shown above.

\ref{prelim:automata:morpisms:reachability:prop:images:twoSCC}: 
This follows from \cref{prelim:general:order:prop:equiv-class-image-disjoint-or-subset}.
\end{proof}

For general order morphisms, \Cref{prelim:eqn:morph-principal-upper-set} may fail and 
the images of upper sets need not be upper sets, see \cref{prelim:general:order:rem:morphism-upperset}.

\begin{corollary} \label{prelim:automata:morpisms:reachability:cor:images}
	Let $T \subseteq Q$.
	If $\varphi(T) \subseteq R$ for a sink component~$R$,
	then $\varphi(\upperset{T}) = R$.
\end{corollary}
\begin{proof}
	Since $\varphi(T) \subseteq R$
	we have $\upperset{\varphi(T)} \subseteq \upperset{R} = R$.
	Since~$R$ is strongly connected and $\upperset{\varphi(T)} \cap R \ne \emptyset$,
	we have $R \subseteq \upperset{\varphi(T)}$.
	Thus, with \cref{prelim:automata:morpisms:reachability:prop:images},
	we find $\varphi(\upperset{T}) = \upperset{\varphi(T)} = R$.
\end{proof}

\begin{remark}
	Consider the automata in \cref{prelim:aut:fig:mappingmorph,prelim:aut:fig:mappingmorph0} 
	from \cref{prelim:ex:SCC2}. 
	The image of $\{q_0\}$ is $\{q_1\}$, which is strictly
	contained in the sink component $\{q_1, q_2\}$,
	and the image of $\upperset{\{q_0\}}$ equals $\{q_1, q_2\}$.
	Note that $\{q_0\}$ is a strongly connected component,
	but its image is not a strongly connected component in both automata.
	
	In \cref{prelim:aut:fig:mappingotmorph}, $\{q_0\}$
	is a sink component and its image is not a sink component.
	Thus, in this case, we also obtain that the mapping
	is not a morphism from \cref{prelim:automata:morpisms:reachability:cor:images}.
\end{remark}

In general, bijective order morphisms
do not have to be order isomorphisms, see \cref{prelim:ex:order:bijmorph}.
By the next statement, with \cref{prelim:prop:iso-bijmorph}, bijective order 
morphisms arising from automata morphisms are in fact order isomorphisms.

\begin{proposition} \label{prelim:automata:morphism:prop:injective}
	If $\varphi \colon Q \to S$ is an isomorphism,
	then it is also an order isomorphism.
\end{proposition}
\begin{proof}
	By assumption, there exists a morphism $\psi \colon S \to Q$
	such that $\concat{\varphi}{\psi} = \operatorname{id}_Q$
	and $\concat{\psi}{\varphi} = \operatorname{id}_S$.
	By \cref{prelim:eqn:morph-reachability-order}, $\psi$
	is also an order morphism.
	Thus, $\varphi$ is an order isomorphism.	
\end{proof}

\begin{corollary} \label{prelim:automata:morphism:prop:injective:SCC-to-SCC}
	Suppose $\varphi \colon Q \to S$ is an isomorphism.
	If $T \subseteq Q$ is a strongly connected component,
	then $\varphi(T)$ is a strongly connected component. 
\end{corollary}
\begin{proof}
	This follows from \cref{prelim:prop:morphism_equiv_class,prelim:automata:morphism:prop:injective}.
\end{proof}


\begin{proposition} \label{prelim:automata:morphism:prop:image-SCC}
\begin{enumerate}
\item \label{prelim:automata:morphism:prop:image-SCC:1} If $Q \ne \emptyset$ and $\mathcal B$ is strongly connected, or strongly trap-connected where $\varphi(Q)$ has size at least two,
then $\varphi$ is surjective, i.e., $\varphi(Q) = S$.

\item \label{prelim:automata:morphism:prop:image-SCC:2} If $\mathcal A$ is strongly connected (resp.~strongly trap-connected)
and $\varphi$ is surjective, 
then~$\mathcal B$ is strongly connected (resp.~strongly trap-connected). 
\end{enumerate}
\end{proposition}
\begin{proof}
\ref{prelim:automata:morphism:prop:image-SCC:1}: By \cref{prelim:automata:morpisms:reachability:prop:images:preservation} of \cref{prelim:automata:morpisms:reachability:prop:images}, $\varphi(Q)$ is an upper set.
 First, suppose $\mathcal B$ is strongly connected.
 By \cref{prelim:automata:prop:char-strongly-conn}, a strongly connected automaton has no proper non-empty upper
 set. Thus $\varphi(Q) = S$.
 Next, suppose~$\mathcal B$ is strongly trap-connected and $|\varphi(Q)| \ge 2$.
 By \cref{prelim:automata:prop:char-strongly-trap}, $\mathcal B$
 has no proper upper set of size at least two.
 Thus $\varphi(Q) = S$.

\ref{prelim:automata:morphism:prop:image-SCC:2}: We can assume $|Q|, |\varphi(Q)| \ge 2$.

If~$\mathcal A$ is strongly connected, then~$Q$ is a sink component.
Thus, by \cref{prelim:automata:morpisms:reachability:prop:images:preservation} from \cref{prelim:automata:morpisms:reachability:prop:images},
$\varphi(Q)$ is also a sink component. By surjectivity, $\mathcal B$ is strongly connected.

Next, suppose~$\mathcal A$ is strongly trap-connected.
Let $T \subseteq S$ be an upper set in $\mathcal B$
of size at least two.
By \cref{prelim:general:order:eqn:morphism-upperset-inverse}, $\varphi^{-1}(T)$
is an upper set. Since~$\varphi$ is surjective, $|\varphi^{-1}(T)| \ge 2$.
By \cref{prelim:automata:prop:char-strongly-trap} we find $Q = \varphi^{-1}(T)$.
Thus $S = \varphi(Q) \subseteq T$.
Since the trap state in $\mathcal A$ is mapped to a trap state in~$\mathcal B$,
$\mathcal B$ is not strongly connected.
Then, \cref{prelim:automata:prop:char-strongly-trap} gives
that $\mathcal B$ is strongly trap-connected.
\end{proof}

\begin{corollary} \label{prelim:automata:morphism:cor:image-SCC}
	If $|Q| = |S|$ and~$\mathcal B$ is strongly connected,
	or strongly trap-connected with $|\varphi(Q)| \ge 2$,
	then~$\varphi$ is an isomorphism (and so~$\mathcal A$ is also strongly connected or strongly trap-connected).
\end{corollary}
\begin{proof}
	By \cref{prelim:automata:morphism:prop:image-SCC:1} from \cref{prelim:automata:morphism:prop:image-SCC}, 
	we have $\varphi(Q) = S$.
	Since $\rank{\varphi} = |Q|$,
	by \cref{prelim:prop:item:rank-defect-eqn}
	from \cref{prelim:prop:rank-defect} we find $\defect{\varphi} = \emptyset$.
	Thus, $\varphi$ is bijective.
	By \cref{prelim:prop:iso-bijmorph}, $\varphi$ is an isomorphism.
	
	The inverse map of~$\varphi$ is an isomorphism.
	Since~$\mathcal B$ is strongly connected or strongly trap-connected,
	by \cref{prelim:automata:morphism:prop:image-SCC:2} from \cref{prelim:automata:morphism:prop:image-SCC},
	we find that $\mathcal A$ is strongly connected
	or strongly trap-connected.		
\end{proof}

\subsection{Endomorphisms}
\label{prelim:sec:endo}

By the first result of this section, \cref{prelim:automata:prop:endo-reach-order},
an endomorphism cannot map a state to one that strictly precedes it in the
reachability order.
This property
does not hold for general quasiorders, see \cref{prelim:general:order:rem:morphism-upperset}.
We derive several consequences of this fact. For example,
\cref{prelim:automata:prop:endo:SCC}, showing that an endomorphism permutes a
strongly connected component~$T$ if and only if it permutes the upper set
generated by~$T$, or
\cref{prelim:automata:cor:SCC-disjoint-or-permutates} showing that an
endomorphism either permutes a strongly connected component or maps all of
its states to states outside the component.
Lastly, we give several results on endomorphisms of strongly connected
and synchronizing automata.

In this section, let $\semiautomat{A}{Q}{\Sigma}{\delta}$ be an automaton and $\varphi \colon Q \to Q$ be an endomorphism. 

\begin{proposition} \label{prelim:automata:prop:endo-reach-order}
 Let $p \in Q$.
 Then $\varphi(p) \lereach p$ implies $p \lereach \varphi(p)$.	
\end{proposition}
\begin{proof}
	Suppose $\varphi(p) \lereach p$.
	By \cref{prelim:eqn:reach-upperset} 
	we have $\upperset{\{p\}}\subseteq \upperset{\{\varphi(p)\}}$.
	By \cref{prelim:eqn:morph-principal-upper-set} 
	we have $\upperset{\{\varphi(p)\}} = \varphi(\upperset{\{p\}})$.
	Thus $\upperset{\{p\}}\subseteq \varphi(\upperset{\{p\}})$.
	By finiteness, this gives 
	$\upperset{\{p\}} = \varphi(\upperset{\{p\}}) = \upperset{\{\varphi(p)\}}$.
	\Cref{prelim:general:order:eqn:equiv-uppersets}
	we find $p \eqreach \varphi(p)$.
\end{proof}

\begin{corollary}
	Either $p$ and $\varphi(p)$ are incomparable or we have $p \lereach \varphi(p)$
\end{corollary}

\begin{corollary} \label{prelim:automata:prop:endo:SCC}
	For a strongly connected component $T \subseteq Q$,
	the following are equivalent:
	\begin{enumerate}
		\item \label{prelim:facts:permutes_SCC:1a} $T \subseteq \varphi(\upperset{T})$,
		
		\item \label{prelim:facts:permutes_SCC:3} $\varphi(T) = T$,

		\item \label{prelim:facts:permutes_SCC:2} $\varphi(T) \cap T \ne \emptyset$,
		
		\item \label{prelim:facts:permutes_SCC:1} $\varphi(\upperset{T}) = \upperset{T}$.
	\end{enumerate}	
\end{corollary}
\begin{proof}

$\ref{prelim:facts:permutes_SCC:1a} \Rightarrow \ref{prelim:facts:permutes_SCC:3}$:
Let $p \in T$. By assumption, we have $p = \varphi(q)$ for some $q \in \upperset{T}$.
Since~$T$ is a strongly connected component,
we have $\upperset{T} = \upperset{\{p\}}$.
Thus, there exists $u \in \Sigma^*$ such that $q = \delta(p, u)$.
This gives $p = \varphi(\delta(p, u))$.
By \cref{prelim:prop:reachability-finest} we have 
$p \lereach \delta(p, u)$.
Hence, $\varphi(\delta(p, u)) \lereach \delta(p, u)$.
By \cref{prelim:automata:prop:endo-reach-order}
this gives $\varphi(\delta(p, u)) \eqreach \delta(p, u)$.
Thus $p \eqreach \delta(p, u)$
and so $q \in T$. Hence $p \in \varphi(T)$.
This gives $T \subseteq \varphi(T)$.
Since $|\varphi(T)| \le |T|$ by \cref{eqn:subset-mapped-to-smaller-set},
we find $T = \varphi(T)$.	

$\ref{prelim:facts:permutes_SCC:3} \Rightarrow \ref{prelim:facts:permutes_SCC:2}$: This is clear.

$\ref{prelim:facts:permutes_SCC:2} \Rightarrow \ref{prelim:facts:permutes_SCC:1}$: Let $p \in \varphi(T) \cap T$.
By \cref{prelim:general:order:prop:equiv-class-image-disjoint-or-subset} 
we have $\varphi(T) \subseteq T$.
Thus $\varphi(p) \in T$.
Since~$T$ is a strongly connected
component, we have $\upperset{T} = \upperset{\{p\}} = \upperset{\{\varphi(p)\}}$. So
\[
\varphi(\upperset{T})
= \varphi(\upperset{\{p\}})
= \upperset{\{\varphi(p)\}}
= \upperset{T}.
\]
This gives the claim.\footnote{Alternatively,
	since $\varphi(\upperset{T})$ is an upper set having non-empty intersection with~$T$,
	we have $T \subseteq \varphi(\upperset{T})$.
	As $\upperset{T}$ is the smallest upper set containing~$T$,
	this gives $\upperset{T}\subseteq \varphi(\upperset{T})$.
	By finiteness, we then obtain 
	$\upperset{T} = \varphi(\upperset{T})$.}
	
$\ref{prelim:facts:permutes_SCC:1} \Rightarrow \ref{prelim:facts:permutes_SCC:1a}$:	This is clear.
\end{proof}

\begin{corollary} \label{prelim:automata:cor:SCC-disjoint-or-permutates}
	For a strongly connected component $T \subseteq Q$,
	we have $\varphi(T) \cap T = \emptyset$
	or $\varphi(T) = T$.
\end{corollary}
\begin{proof}
	This directly follows from Items~\ref{prelim:facts:permutes_SCC:3} and~\ref{prelim:facts:permutes_SCC:2}
	of the previous corollary.	
\end{proof}

\begin{remark}	
	In general, for a sink component $T \subseteq Q$, a morphism $\varphi$ need not induce an isomorphism between
	$\mathcal A_{|T}$ and $\mathcal A_{|\varphi(T)}$.
	For example, this fails when $|\varphi(T)| = 1$ and $|T| > 1$.
	However, by \cref{prelim:automata:cor:SCC-disjoint-or-permutates}, if
	$\varphi(T) \cap T \neq \emptyset$, then $\varphi(T)=T$.
	Consequently, the restriction $\varphi_{|T}$ is an automorphism of
	$\mathcal A_{|T}$.
\end{remark}

\begin{corollary} \label{prelim:automata:morphism:prop:endo-SCC-is-iso}
	An endomorphism of a strongly connected finite automaton is an automorphism.
\end{corollary}
\begin{proof}
	By \cref{prelim:automata:cor:SCC-disjoint-or-permutates},
	it follows that every endomorphism is bijective, and so an automorphism by \cref{prelim:prop:iso-bijmorph}.  
	
	Alternatively, this also follows from \cref{prelim:automata:morphism:prop:image-SCC:1} 
	of \cref{prelim:automata:morphism:prop:image-SCC}.
\end{proof}

There exist automata that are not strongly connected and whose non-identity endomorphisms are all automorphisms. For example, a unary automaton in which the single letter induces two cycles of coprime lengths.

\begin{corollary} \label{prelim:automata:morphim:cor:strongly-conn}
	Let $S \subseteq Q$ be strongly connected.
	Then $\varphi(S) \subseteq S$
	implies $\varphi(S) = S$.
\end{corollary}
\begin{proof}
	We have $S \subseteq T$ for some strongly connected component $T \subseteq Q$.
	Let $p \in S$. By assumption, we have $\varphi(p) \in T$.
	Thus $\varphi(p) \in T \cap \varphi(T)$.
	By \cref{prelim:automata:cor:SCC-disjoint-or-permutates}
	we find $\varphi(T) = T$.
	Hence~$\varphi$ permutes the states in~$T$
	and so $\varphi_{|T}$ is injective.
	Then $\varphi(S) \subseteq S$ gives $\varphi(S) = S$.
\end{proof}

\begin{corollary}
	Let $\psi \colon Q \to Q$ be another endomorphism
	and $S \subseteq Q$ be strongly connected.
	If $\psi(\varphi(S)) \subseteq \varphi(S)$,
	then $\psi(\varphi(S)) = \varphi(S)$.	
\end{corollary}
\begin{proof}
	By \cref{prelim:automata:morpisms:reachability:prop:images:twoSCC} of \cref{prelim:automata:morpisms:reachability:prop:images}
	the set $\varphi(S)$ is strongly connected.
	Then the claim follows from \cref{prelim:automata:morphim:cor:strongly-conn}.	
\end{proof}

\begin{corollary}
	For $\varphi(S) \subseteq S$ and $\delta(S, u) \subseteq T$ with strongly connected components~$S, T \subseteq Q$ 
	and $u \in \Sigma^*$ we have $\varphi(T) = T$.	
\end{corollary}
\begin{proof}
	Let $q \in S$. Then $\varphi(q) \in S$.
	Since $\delta(S, u) \subseteq T$,
	we find $\delta(q, u) \in T$ and $\varphi(\delta(q, u)) = \delta(\varphi(q), u) \in T$.	
	Thus $\varphi(\delta(q, u)) \in \varphi(T) \cap T$.
	By \cref{prelim:automata:prop:endo:SCC},
	this gives $\varphi(T) = T$.
\end{proof}

By \cref{prelim:automata:prop:endo:SCC},
for $\varphi(\upperset{\{p\}}) = \upperset{\{p\}}$, $\varphi$ permutes the strongly connected component
$[p]_{\eqreach}$.
However, this does not imply
that $\varphi$ permutes each strongly connected component in $\upperset{\{p\}}$,
as \cref{prelim:ex:endo} shows next.

\begin{example} \label{prelim:ex:endo}
	Let $\semiautomat{A}{Q}{\{a,b\}}{\delta}$
	with $Q = \{q_1, q_2, q_3, q_4\}$ and
	\[
	\delta(q_i, a) = 
	\begin{cases}
		q_4 & \text{if } i = 1, \\
		q_3 & \text{if } i = 2, \\
		q_i & \text{otherwise,} 
	\end{cases}
	\qquad 
	\delta(q_i, b) =
	\begin{cases}
		q_1  & \text{if } i = 2, \\
		q_2  & \text{if } i = 1, \\
		q_i  & \text{otherwise.} 
	\end{cases}
	\]
	The strongly connected components are $\{q_1, q_2\}$, $\{q_3\}$ and $\{q_4\}$.
	We have $\upperset{\{q_1\}} = Q$.
	The map $\varphi$ given by $\varphi(q_1) = q_2$, $\varphi(q_2) = q_1$, $\varphi(q_3) = q_4$
	and $\varphi(q_4) = q_3$ is a morphism
	that maps $\{q_1, q_2\}$ to itself and interchanges $\{q_3\}$ and $\{q_4\}$.
	
	\begin{figure}[ht]
		\centering	
		\begin{tikzpicture}[node distance = 2cm, on grid, auto, ->]
			\node (q1) [state] {$q_1$};
			\node (q2) [state, right of = q1] {$q_2$};
			\node (q3) [state, right of = q2] {$q_3$};
			\node (q4) [state, left of  = q1] {$q_4$};
			
			\path (q1) edge [above] node {$a$} (q4)
			      (q2) edge node {$a$} (q3)
			      (q3) edge [loop right] node {$a,b$} (q3)
			      (q4) edge [loop left]  node {$a,b$} (q4);
			\path (q1) edge [bend left] node {$b$} (q2)
			      (q2) edge [bend left] node {$b$} (q1); 
		\end{tikzpicture}
		\caption[Some strongly connected component fixed and others swapped]{Automaton admitting a morphism that leaves $\{q_1, q_2\}$ invariant and swaps $\{q_4\}$ and $\{q_3\}$, see \cref{prelim:ex:endo}.}
		\label{prelim:aut:fig:autnomorph2}
	\end{figure}
\end{example} 

\begin{proposition} \label{prelim:automata:prop:endo-SCC-fixed-all-or-none}
	For a strongly connected component $T \subseteq Q$, 
	either we have $\varphi(p) = p$ for all $p \in T$
	or $\varphi(p) \ne p$ for all $p \in T$.
\end{proposition}
\begin{proof}
	Suppose we have $\varphi(p) = p$ for some $p \in T$.
	Since $T \subseteq \upperset{\{p\}}$,
	from \cref{prelim:automata:morphism:cor:two-morp-upper-set} 
	applied to~$\varphi$ 
	and the identity transformation~$\operatorname{id}_Q$
	we get $\varphi(q) = q$ for all $q\in \upperset\{p\}$,
	in particular~$\varphi$ fixes all states of~$T$.	
\end{proof}

\begin{corollary} \label{prelim:automata:morphism:prop:endo-SCC-is-iso-b}
    The automorphism group of a strongly connected automaton is semi-regular (and so has size at most~$|Q|$).
\end{corollary}
\begin{proof}
	By  \cref{prelim:automata:prop:endo-SCC-fixed-all-or-none}, every non-identity isomorphism
	fixes no state.
	Thus, the automorphism group is semiregular.
	By \cref{prelim:groups:prop:semiregular-size}, 
	its size is at most~$|Q|$.	
\end{proof}

\begin{proposition} \label{prelim:automata:prop:endo-SCC-fixed-sync-state}
	$\varphi \colon Q \to Q$ maps a synchronizing state to itself.
\end{proposition}
\begin{proof}
	Let $u \in \Sigma^*$ be such that $\delta(Q, u) = \{q\}$.
	Then we have $\varphi(q) = \varphi(\delta(q, u)) = \delta(\varphi(q), u) = q$.
\end{proof}

\begin{corollary} \label{prelim:automata:morphism:cor:endo-sync-SCC}
	The only endomorphism of a strongly connected synchronizing automaton
	is the identity transformation.
\end{corollary}
\begin{proof}
	By \cref{prelim:automata:prop:endo-SCC-fixed-sync-state}
	an endomorphism fixes the synchronizing state.
	Thus, by \cref{prelim:automata:prop:endo-SCC-fixed-all-or-none},
	it equals the identity transformation.\footnote{Alternatively,
	since in this case every state is a synchronizing state
	as we show later in \cref{prelim:cor:SCC_sync},
	from \cref{prelim:automata:prop:endo-SCC-fixed-sync-state}
	we also get that the endomorphism fixes every state.}		
\end{proof}

\begin{corollary}
	For a strongly connected synchronizing automaton
	having the same number of states as~$\mathcal A$,
	there exists at most one morphism from~$\mathcal A$ to that automaton.
\end{corollary}
\begin{proof}
	Let $\semiautomat{B}{S}{\Sigma}{\mu}$ be strongly connected
	and synchronizing with $|Q| = |S|$.
	Suppose $\psi, \xi \colon Q \to S$ are two morphisms.
	Since~$\mathcal B$ is strongly connected,
	by \cref{prelim:automata:morphism:cor:image-SCC}, 
	they are both isomorphisms.
	Let $\psi^{-1} \colon S \to Q$ be the inverse morphism of~$\psi$.
	Then $\concat{\psi^{-1}}{\xi}$ is an endomorphism of~$\mathcal B$.
	By \cref{prelim:automata:morphism:cor:endo-sync-SCC} 
	it equals the identity transformation, which implies $\psi = \xi$.
\end{proof}

\subsection{Cyclic Unary Automata}

As a simple example, here we classify all morphisms
between cyclic unary automata. First, \cref{prelim:automata:lem:morphism-unary-automata,prelim:automata:prop:morphism-unary-automata}
are preliminary statements, and \cref{prelim:automata:cor:morphism-unary-automata} 
gives the classification.

\begin{lemma} \label{prelim:automata:lem:morphism-unary-automata}
	Let $\semiautomat{A}{Q}{\{a\}}{\delta}$
	and $\semiautomat{B}{S}{\{a\}}{\mu}$
	where $Q = \upperset{\{p\}}$ for some $p \in Q$.
	Suppose $\varphi \colon Q \to S$ is a morphism.
	Let $i_2 \ge 0$ be the least number such that,
	with a corresponding least number $p_2 > 0$,
	we have $\mu(\delta(p), a^{i_2 + p_2}) = \mu(\delta(p), a^{i_2})$.
	Then for all  $i_1 \ge 0$ and $p_1 > 0$
	with $\delta(p, a^{i_1 + p_1}) = \delta(p, a^{i_1})$ 
	we have  $i_2 \le i_1$ and that $p_1$ is a multiple of~$p_2$. 
\end{lemma}
\begin{proof}
	First, $\delta(p, a^{i_1 + p_1}) = \delta(p, a^{i_1})$
	implies $\mu(\varphi(p), a^{i_1 + p_1}) = 
	\varphi(\delta(p, a^{i_1 + p_1})) = \varphi(\delta(p, a^{i_1})) = \mu(\varphi(p), a^{i_1})$.
	Hence, by the minimal choice of~$i_2$ and~$p_2$,
	we have $i_2 \le i_1$ and $p_2 \le p_1$.
	Next, write $p_1 = l p_2 + r$
	for some $l \ge 0$ and $0 \le r < p_2$.
	Then 
	we have $\mu(\delta(p), a^{i_2 + p_1}) = \mu(\delta(p), a^{i_2 + r})$
	and, by minimality, we must have $r = 0$, i.e., $p_1$
	is a multiple of~$p_2$.	
\end{proof}

\begin{example}
	Let $\mathcal A = (Q, \{a\}, \delta)$ consist of two disjoint
	cycles for~$a$ of lengths $p$ and~$q$
	for two distinct prime numbers.
	Then endomorphisms mapping some state from one cycle
	to the other cycle are not possible by \cref{prelim:automata:lem:morphism-unary-automata}.
	With \cref{prelim:automata:cor:endo-unary-automata}, each endomorphism
	shifts each state on the cycle of length~$p$
	and on the cycle of length~$q$ by some
	amount~$0 \le i < p$ and $0 \le j < q$, respectively.
	So, $\mathcal A$
	has~$pq$ endomorphisms.
\end{example}

\begin{corollary} \label{prelim:automata:cor:morphism-unary-automata}
	Let $\semiautomat{A}{Q}{\{a\}}{\delta}$ 
	where $Q = \upperset{\{p\}}$ for some $p \in Q$
	and $i_1 \ge 0$ be the least number
	with a corresponding least number $p_1 > 0$
	such that $\delta(p, a^{i_1+p_1}) = \delta(p, a^{i_1})$.
	Then $\delta(p, a^{i_2 + p_2}) = \delta(p, a^{i_2})$
	implies $i_1 \le i_2$ and that $p_2$ is a multiple of~$p_1$.
\end{corollary}
\begin{proof}
	Apply \cref{prelim:automata:lem:morphism-unary-automata} with the identity transformation, which is a morphism.
\end{proof}

\begin{proposition} \label{prelim:automata:prop:morphism-unary-automata}
	Let $\semiautomat{A}{Q}{\{a\}}{\delta}$
	with $Q = \upperset{\{q\}}$ for some $q \in Q$
	and $\semiautomat{B}{S}{\{a\}}{\mu}$.
	Then $\varphi \colon Q \to S$ is a morphism if and only if $\mu(\varphi(q), a^{i_2 + p_2}) = \mu(\varphi(q), a^{i_2})$
	for some $0 \le i_2 \le i_1$ with $p_1$ being a multiple of~$p_2$
	and $\varphi(\delta(q, a^i)) = \mu(\varphi(q), a^i)$ for $i \in\{0,1,\ldots,n-1\}$
	with $n = |Q|$.
\end{proposition}
\begin{proof}
	First, let $\varphi$ be a morphism.
	By finiteness, we have  $\mu(\varphi(q), a^{i_2 + p_2}) = \mu(\varphi(q), a^{i_2})$
	for some minimal $i_2 \ge 0$ and some $p_2 > 0$.
	By \cref{prelim:automata:lem:morphism-unary-automata},
	we find $i_2 \le i_1$ and that $p_1$ is a multiple of~$p_2$.
	Since $\varphi$ is a morphism,
	we have $\varphi(\delta(q, a^i)) = \mu(\varphi(q), a^i)$ for all $i \ge 0$.
	
	Next, suppose we have $\mu(\varphi(q), a^{i_2 + p_2}) = \mu(\varphi(q), a^{i_2})$
	for some $0 \le i_2 \le i_1$ with $p_1$ being a multiple of~$p_2$
	and  $\varphi(\delta(q, a^i)) = \mu(\varphi(q), a^i)$ for $i \in\{0,1,\ldots,n-1\}$.
	We have $n = i_1 + p_1$.
	By assumption, there exists $k \ge 1$ such that $p_1 = k\cdot p_2$.
	Since $i_2 \le i_1$ 
	we have $\mu(\varphi(q), a^{i_1}) = \mu(\varphi(q), a^{i_1 + k\cdot p_2}) = \mu(\varphi(q), a^n)$.
	Then we find	
	\[
	\varphi(\delta(q, a^{n - 1}), a) = \varphi(\delta(q, a^{i_1}))
	= \mu(\varphi(q), a^{i_1})
	= \mu(\varphi(q), a^n)
	= \mu(\delta(q, a^{n - 1}), a).
	\]
	Since $Q = \{q, \delta(q, a), \ldots, \delta(q, a^{n-1})\}$, 
	we have $\varphi(\delta(p, a)) = \delta(\varphi(p), a)$
	for all $p \in Q$.
\end{proof}

\begin{corollary} \label{prelim:automata:cor:morphism-unary-automata}
	Let $\semiautomat{A}{Q}{\{a\}}{\delta}$
	and $\semiautomat{B}{S}{\{a\}}{\mu}$
	where $Q = \upperset{\{p\}}$ 
	and $S = \upperset{\{s\}}$
	for some $q \in Q$ and $s \in S$.
	Let  $i_1, i_2 \ge 0$ be the least numbers with corresponding least numbers $p_1,p_2 > 0$ 
	such that $\delta(q, a^{i_1 + p_1}) = \delta(q, a^{i_1})$
	and $\mu(s, a^{i_2 + p_2}) = \mu(s, a^{i_2})$.
	Then there exist $\min\{i_1, i_2\} + p_2$
	morphisms from~$\mathcal A$ to~$\mathcal B$ 
	if and only if $p_1$ is a multiple of~$p_2$,
	and there exists no morphism otherwise.
\end{corollary}
\begin{proof}
	If $p_1$ is not a multiple of~$p_2$,
	there exists no morphism by \cref{prelim:automata:cor:morphism-unary-automata}.
	Next, suppose $p_1$ is a multiple of~$p_2$.
	Then, for $i_2 - \min\{i_1, i_2\} \le i < |S|$,
	setting $\varphi(\delta(q, a^j)) = \mu(s, a^{i + j})$,
	which gives  $\varphi(\delta(q, a^j)) = \mu(\varphi(q), a^j)$,
	for $j \in \{0,1,\ldots,|Q|-1\}$
	yields a morphism by \cref{prelim:automata:cor:morphism-unary-automata}
	where $\mu(\varphi(q), a^{\min\{i_2 - i, 0\} + p_2}) = \mu(\varphi(q), a^{\min\{i_2 - i,0\}})$
	and $\min\{i_2 - i, 0\} \le \min\{i_1, i_2\} \le i_1$.
	
	Next, let $\varphi \colon Q \to S$ be a morphism
	and $0 \le i <|S|$ such that $\varphi(q) = \mu(s, a^i)$.
	We have $\mu(s, a^{i + j}) = \mu(s, a^{i + j + p_2})$
	for $j = \min\{i_2 - i, 0\}$.
	By \cref{prelim:automata:cor:morphism-unary-automata} we find $j \le i_1$.
	Thus, for $i_2 > i_1$ we must have $i_2 - i_1 \le i$
	since otherwise $i + j < i_2$.
	This gives $i_2 - \min\{i_1, i_2\} \le i < |S|$.
	
	So, each morphism is uniquely
	given by its image 
	$\varphi(q) = \mu(s, a^i)$
	for $i_2 - \min\{i_1, i_2\} \le i < |S|$.
	Thus for $i_2 < i_1$
	we have $|S| = i_2 + p_2 = \min\{i_1, i_2\} + p_2$ 
	possible morphisms.
	For $i_2 \ge i_1$
	we have $|\{ i \mid i_2 - i_1 \le i < i_2 + p_2 \}| =i_1 + p_2 = \min\{i_1, i_2\} + p_2$
	possible morphisms.	
\end{proof}

\begin{corollary} \label{prelim:automata:cor:endo-unary-automata}
	Let $\semiautomat{A}{Q}{\{a\}}{\delta}$
	with $Q = \upperset{\{q\}}$ for some $q \in Q$.
	Then $\varphi \colon Q \to Q$ is an endomorphism
	if and only if there exists $0 \le m < |Q|$ such that
	$\varphi(p) = \delta(p, a^m)$
	for all $p \in Q$.
\end{corollary}
\begin{proof}
	Let $m \ge 0$ such that $\varphi(q) = \delta(q, a^m)$.
	By \cref{prelim:automata:cor:morphism-unary-automata}, $\varphi$
	is a morphism if and only if $\varphi(\delta(q, a^i)) = \delta(\varphi(q), a^i)$
	for $0 \le i < |Q|$.
	This is equivalent to $\varphi(p) = \delta(p, a^m)$
	for all $p \in Q$ since $\delta(\varphi(q), a^i) = \delta(\delta(q, a^i), a^m)$.
\end{proof}

Since in unary automata, strongly connected components are cycles, by \cref{prelim:automata:lem:morphism-unary-automata}, 
the size of a strongly connected component in the image of a morphism divides the size of a strongly connected
component that maps to it. This might fail for non-unary automata, as shown next.

\begin{example} \label{prelim:aut:ex:morphism-not-divide}
	Let $\semiautomat{A}{\{q_0, q_1, q_2\}}{\{a,b\}}{\delta}$
	and
	$\semiautomat{B}{\{p_0, p_1\}}{\{a,b\}}{\mu}$
	with
	$\delta(q_0, a) = q_1$, $\delta(q_1, a) = q_2$,
	$\delta(q_1, b) = q_0$, $\delta(q_2, b) = q_0$,
	$\mu(p_0, a) = p_1$, $\delta(p_1, b) = p_0$
	and self-loops for all the remaining transitions, see \cref{prelim:aut:fig:morphism-not-divide}.
	The map~$\varphi$ with $\varphi(q_0) = p_0$, $\varphi(q_1) = p_1$,
	$\varphi(q_2) = p_1$
	is a morphism since $\mu(\varphi(q), b) = p_0 = \varphi(q_0) = \varphi(\delta(q, a))$
	for each $q \in \{q_0,q_1,q_2\}$ and
	\begin{align*}
		\mu(\varphi(q_0), a) & = \mu(p_0, a) = p_1 = \varphi(q_1) = \varphi(\delta(q_0, a)), \\
		\mu(\varphi(q_1), a) & = \mu(p_1, a) = p_1 = \varphi(q_2) = \varphi(\delta(q_1, a)), \\
		\mu(\varphi(q_2), a) & = \mu(p_1, a) = p_1 = \varphi(q_2) = \varphi(\delta(q_2, a)).
	\end{align*}	
	
	\begin{figure}[ht]
		\centering	
		\begin{tikzpicture}[->, >=Stealth, node distance=8em, on grid, auto]
			
			% --- Automaton A ---
			\node[state] (q0) {$q_0$};
			\node[state, right=of q0] (q1) {$q_1$};
			\node[state, right=of q1] (q2) {$q_2$};
			
			% Transitions in A
			\path
			(q0) edge[bend left=10] node {$a$} (q1)
			(q1) edge[bend left=10] node {$b$} (q0)
			(q1) edge[bend left=10] node {$a$} (q2)
			(q2) edge[bend right=32, above] node {$b$} (q0);
			
			%	% Self-loops in A
			%	(q0) edge[loop above] node {$b$} ()
			%	(q2) edge[loop above] node {$a$} ();
			
			% Label A
			\node[right=14em of q1] {Automaton $\mathcal A$};
			
			% --- Automaton B ---
			\node[state, below=7em of q0] (p0) {$p_0$};
			\node[state, right=of p0] (p1) {$p_1$};
			
			% Transitions in B
			\path
			(p0) edge[bend left=10] node {$a$} (p1)
			(p1) edge[bend left=10] node {$b$} (p0);
			
			%	% Self-loops in B
			%	\path
			%	(p0) edge[loop below] node {$b$} ()
			%	(p1) edge[loop below] node {$a$} ();
			
			% Label B
			\node[right=6em of p1] {Automaton $\mathcal B$};
			
			% --- Morphism phi ---
			\draw[dashed]
			(q0) -- node[left=0.6cm] {$\varphi$} (p0);
			\draw[dashed] (q1) -- (p1);
			\draw[dashed] (q2) -- (p1);
			
		\end{tikzpicture}
		\caption[Morphism between automata]{Morphism $\varphi$ between two automata $\mathcal A$ and~$\mathcal B$, see \cref{prelim:aut:ex:morphism-not-divide}.}
		\label{prelim:aut:fig:morphism-not-divide}
	\end{figure}		
\end{example}